% 20260906 Science Quiz mathematics Problem set.

% Verified against the official Abel Prize 2025 announcement and citation.
\begin{problem}{1}
$D$-modules 이론의 발전과 crystal bases의 발견으로 잘 알려진 Masaki Kashiwara가 2025년에 수상한 주요 국제 수학상을 쓰시오.
\end{problem}

\begin{solution}
The Norwegian Academy of Science and Letters는 Masaki Kashiwara에게 algebraic analysis와 representation theory에 대한 근본적 기여, 특히 $D$-modules 이론의 발전과 crystal bases의 발견을 이유로 2025 Abel Prize를 수여했다.

따라서 정답은
\[
    \boxed{\text{Abel Prize}}
\]
이다.
\end{solution}

\begin{problem}[Broken Stick Problem]{2}
길이 $1$인 선분 위에서 서로 독립적으로 uniformly random한 두 점을 골라 선분을 세 조각으로 자른다. 세 조각이 삼각형의 세 변의 길이가 될 확률을 구하시오.
\end{problem}

\begin{solution}
두 절단점의 위치를 $U,V\in(0,1)$이라 하고, 대칭성에 의해 $0<U<V<1$인 경우를 먼저 생각하자. 세 조각의 길이는
\[
    U,
    \qquad
    V-U,
    \qquad
    1-V
\]
이다. 세 양수가 삼각형의 세 변이 될 필요충분조건은 가장 긴 조각이 $1/2$보다 작은 것이므로
\[
    U<\frac12,
    \qquad
    V-U<\frac12,
    \qquad
    1-V<\frac12.
\]
즉
\[
    0<U<\frac12,
    \qquad
    \frac12<V<U+\frac12.
\]
단위 정사각형의 절반인 영역 $U<V$ 안에서 이 조건을 만족하는 영역의 넓이는
\[
    \int_0^{1/2}U\,\dd U
    =\frac18.
\]
$V<U$인 대칭적인 경우까지 합치면 확률은
\[
    2\cdot\frac18
    =\boxed{\frac14}.
\]
\end{solution}

\begin{problem}{3}
함수 $y$가 initial value problem
\[
    \frac{\dd y}{\dd x}=y(1-y),
    \qquad
    y(0)=\frac12
\]
을 만족한다. $y(\ln 3)$을 구하시오.
\end{problem}

\begin{solution}
변수를 분리하면
\[
    \frac{1}{y(1-y)}\,\dd y
    =\dd x.
\]
적분하여
\[
    \ln\frac{y}{1-y}=x+c
\]
를 얻는다. $y(0)=1/2$이므로 $c=0$이고,
\[
    \frac{y}{1-y}=e^x.
\]
따라서
\[
    y(x)=\frac{e^x}{1+e^x}.
\]
$x=\ln3$을 대입하면
\[
    \boxed{y(\ln3)=\frac34}.
\]
\end{solution}

% Verified against the official 2026 Shaw Prize announcement and citation.
\begin{problem}{4}
2026 Shaw Prize in Mathematical Sciences는 mathematical analysis의 깊은 기법을 이용하여 한쪽에서는 information theory, signal processing, statistics의 응용 문제를, 다른 한쪽에서는 geometric measure theory와 fluid dynamics의 singularity를 연구한 두 수학자에게 공동 수여되었다. 두 수상자의 이름을 쓰시오.
\end{problem}

\begin{solution}
2026 Shaw Prize in Mathematical Sciences의 공동 수상자는 Emmanuel Cand\`es와 Camillo De Lellis이다. 공식 citation은 Cand\`es 쪽의 information theory, signal processing, statistics와 De Lellis 쪽의 geometric measure theory, fluid dynamics를 mathematical analysis의 깊은 기법으로 연결한 업적을 강조한다.

따라서 정답은
\[
    \boxed{\text{Emmanuel Cand\`es and Camillo De Lellis}}
\]
이다.
\end{solution}

\begin{problem}{5}
$\R^3$의 두 벡터
\[
    \mathbf{u}\coloneqq(1,1,0),
    \qquad
    \mathbf{v}\coloneqq(1,0,1)
\]
과 원점이 이루는 삼각형의 넓이를 구하시오.
\end{problem}

\begin{solution}
삼각형의 넓이는 두 변 벡터의 cross product 크기의 절반이다. 여기서
\[
    \mathbf{u}\times\mathbf{v}
    =
    \begin{pmatrix}
        1\\-1\\-1
    \end{pmatrix},
\]
이므로
\[
    \|\mathbf{u}\times\mathbf{v}\|
    =\sqrt3.
\]
따라서 넓이는
\[
    \boxed{\frac{\sqrt3}{2}}.
\]
\end{solution}

\begin{problem}{6}
다음을 만족하는 residue class $x\in\Z/105\Z$의 개수를 구하시오.
\[
    x^2\equiv1\pmod{105}.
\]
\end{problem}

\begin{solution}
\[
    105=3\cdot5\cdot7
\]
이고 세 소수는 pairwise coprime이다. 각 소수 $p\in\{3,5,7\}$에 대하여
\[
    x^2\equiv1\pmod p
\]
의 해는
\[
    x\equiv1\pmod p
    \qquad\text{또는}\qquad
    x\equiv-1\pmod p
\]
의 두 개이다. Chinese remainder theorem에 의해 세 modulus에서의 부호 선택 하나마다 modulo $105$의 residue class가 정확히 하나씩 대응한다.

따라서 해의 개수는
\[
    2^3=\boxed{8}.
\]
\end{solution}

% Verified against Alpöge--Furman, arXiv:2608.13637, first posted in August 2026.
\begin{problem}{7}
2026년 8월 Levent Alp\"oge와 Ralph Furman이 공개한 preprint는 Riemann zeta function의 nontrivial zeros 가운데, multiplicity를 포함하여 셌을 때 적어도 얼마의 비율이 simple이면서 critical line 위에 있음을 별도의 가정 없이 보였는가? 분수로 답하시오.
\end{problem}

\begin{solution}
Alp\"oge와 Furman의 2026년 8월 preprint는 Riemann zeta function의 nontrivial zeros를 multiplicity를 포함하여 셌을 때, 적어도 $2/3$가 simple이면서 critical line 위에 있음을 별도의 가정 없이 보인다고 명시한다.

따라서 정답은
\[
    \boxed{\frac23}.
\]
\end{solution}

\begin{problem}{8}
다음 극한을 구하시오.
\[
    \lim_{x\to0}
    \frac{\tan x-\sin x}{x^3}.
\]
\end{problem}

\begin{solution}
$x=0$ 근방에서 Taylor expansion을 사용하면
\[
    \tan x
    =x+\frac{x^3}{3}+O(x^5),
    \qquad
    \sin x
    =x-\frac{x^3}{6}+O(x^5).
\]
따라서
\[
    \tan x-\sin x
    =\frac{x^3}{2}+O(x^5),
\]
이므로
\[
    \boxed{
        \lim_{x\to0}
        \frac{\tan x-\sin x}{x^3}
        =\frac12
    }.
\]
\end{solution}

\begin{problem}{9}
유한 tree의 모든 vertex degree가 $1$, $2$, $3$ 중 하나이고, degree가 $3$인 vertex가 정확히 $4$개이다. 이 tree의 leaf 개수를 구하시오.
\end{problem}

\begin{solution}
Degree가 $1$, $2$, $3$인 vertex의 개수를 각각 $n_1,n_2,n_3$라 하자. 주어진 조건에 의해 $n_3=4$이다. Tree의 vertex 수를 $n$이라 하면
\[
    n=n_1+n_2+n_3
\]
이고 edge 수는 $n-1$이다. Handshaking lemma에 의해
\[
    n_1+2n_2+3n_3
    =2(n-1)
    =2(n_1+n_2+n_3-1).
\]
정리하면
\[
    n_1=n_3+2=6.
\]
Leaf는 degree가 $1$인 vertex이므로 정답은
\[
    \boxed{6}.
\]
\end{solution}

\begin{problem}{10}
정사각형의 네 꼭짓점을 세 가지 색 중 하나로 각각 칠한다. 회전으로 서로 겹치는 coloring은 같은 것으로 보고, reflection으로만 겹치는 coloring은 다른 것으로 본다. 서로 다른 coloring의 개수를 구하시오.
\end{problem}

\begin{solution}
네 번의 회전에 대하여 고정되는 coloring의 수를 센다.

Identity rotation은 모든 coloring $3^4=81$개를 고정한다. $90^\circ$와 $270^\circ$ 회전에 고정되려면 네 꼭짓점이 모두 같은 색이어야 하므로 각각 $3$개이다. $180^\circ$ 회전에 고정되려면 서로 마주보는 두 꼭짓점끼리 같은 색이면 되므로 $3^2=9$개이다.

Burnside's lemma에 의해 서로 다른 coloring의 수는
\[
    \frac{81+3+9+3}{4}
    =\boxed{24}.
\]
\end{solution}
